%%
%% Ergänzung zu Kapitel 4: native iOS-/iPadOS-Plattform
%%

\begin{neu}
\section{Native iOS- und iPadOS-Plattformunterstützung}

Neben der Android-Anwendung wurde eine eigenständige native Portierung für iOS und iPadOS implementiert. Die Portierung ist kein Web-Wrapper und verwendet keine plattformübergreifende UI-Laufzeit, sondern ein eigenes Xcode-Projekt mit Swift und SwiftUI. Android- und Apple-Client besitzen damit getrennte plattformspezifische Implementierungen, orientieren sich jedoch an denselben fachlichen Studienregeln und Backendverträgen. Diese Trennung vermeidet zusätzliche Laufzeitabhängigkeiten und erlaubt es, Sicherheits-, Persistenz- und Hintergrundmechanismen jeweils mit den nativen Betriebssystemdiensten umzusetzen.

Das Deployment Target der Apple-Anwendung liegt bei iOS beziehungsweise iPadOS 17.0. Die Implementierung wurde so angelegt, dass dasselbe App-Target sowohl auf iPhone als auch auf iPad betrieben werden kann. Für Entwicklung und Qualitätssicherung existieren getrennte Konfigurationen für \texttt{Debug}, \texttt{Staging} und \texttt{Release}. Debug verwendet bewusst nicht routbare Testendpunkte, Staging ein freigegebenes lokales Testsystem und Release ausschließlich produktive HTTPS-Endpunkte. Die Konfiguration ist damit bereits auf Ebene des Buildsystems getrennt und nicht von einer zur Laufzeit auswählbaren Serveradresse abhängig.

Die Studienressourcen wurden nicht für die Apple-Plattform neu erzeugt. Stattdessen enthält das iOS-Projekt bytegleiche Kopien der 50 Cue-, 50 Match-A- und 50 Match-B-Bilder der Android-Studienanwendung. Ein Assetmanifest dokumentiert Hashwerte und Abmessungen; ein zusätzliches Contentmanifest beschreibt die Zuordnung von Cue-, Matching- und Labeling-Inhalten. Generator- und Prüfscripte kontrollieren, dass die Kopien deterministisch aus den festgelegten Quellen hervorgehen. Dadurch wird die Plattformportierung von einer inhaltlichen Änderung der experimentellen Stimuli getrennt.

Für die Apple-Plattform werden dieselben zwanzig produktiven Studiensituationen abgebildet. Situationen 1 bis 10 verwenden Cue-Matching, Situationen 11 bis 20 Cue-Labeling. Jeder Durchlauf umfasst fünf Trials und endet mit einer ganzzahligen Rauchverlangensangabe von 0 bis 100 mit dem initialen Sliderwert 50. Bei Matching-Trials bleibt eine Auswahl während der ersten vier sichtbaren Vordergrundsekunden gesperrt. Die 50 Matching-Items werden einmalig in eine zufällige Reihenfolge gebracht und in zehn Blöcke zu je fünf Elementen aufgeteilt; die Labeling-Items folgen den festgelegten Blöcken. Die Position der angebotenen Auswahlmöglichkeiten wird pro Trial neu zufällig festgelegt. Konkrete Trialentscheidungen und Reaktionszeiten werden weder dauerhaft gespeichert noch an das Backend übertragen.

Die plattformspezifischen Betriebssystemfunktionen unterscheiden sich von Android, ohne die fachlichen Studienregeln zu verändern. Das App-Token wird unter iOS im Keychain gespeichert, der Studienfortschritt in einer atomar geschriebenen und vom Backup ausgeschlossenen Datei im App-Container. Lokale Erinnerungen verwenden \texttt{UserNotifications}; eine best-effort Hintergrundprüfung des Informationsfeeds verwendet \texttt{BGAppRefreshTask}. Da iOS den tatsächlichen Ausführungszeitpunkt einer solchen Hintergrundaufgabe selbst bestimmt, bleibt der Feedabruf beim Aktivieren der App der autoritative Weg. Die Kernfunktion der Studie hängt somit nicht von einer garantierten Hintergrundausführung ab.

Die Portierung erweitert die technisch unterstützte Geräteklasse, ohne automatisch den methodischen Geltungsbereich bereits erhobener Daten zu verändern. Ob iPhone- oder iPad-Nutzende tatsächlich in die Wirksamkeitsauswertung eingehen können, hängt zusätzlich davon ab, ob die Apple-Version im freigegebenen Erhebungszeitraum bereitgestellt und im Studienprotokoll beziehungsweise in der Teilnehmendeninformation abgedeckt war. Eine nachträglich implementierte Plattform darf daher nicht rückwirkend als Bestandteil einer zuvor ausschließlich mit Android durchgeführten Erhebung dargestellt werden.

Zum dokumentierten Stand liegt eine vollständige implementierte iOS-/iPadOS-Softwarebasis einschließlich produktivem Studienablauf, Aktivierung, Feed, Feedback, Erinnerungen, Submission, Recovery und beider Abschlussmodi vor. Davon zu unterscheiden ist die formale Distribution: Das vorhandene Audit bewertet die automatisierbaren Implementierungsbereiche als erfüllt, hält die Veröffentlichung jedoch weiterhin zurück, solange unter anderem Distribution-Signing, reale Staging-Ende-zu-Ende-Prüfungen, TestFlight-Gerätetests und externe Freigaben noch offen sind. Die korrekte Formulierung für den Entwicklungsstand lautet deshalb \emph{vollständig implementierter Source Candidate}, nicht \emph{veröffentlichte App-Store-Version}.
\end{neu}
